\section{Zugtiefe}
Die Zugtiefe ist die Anzahl an Schachzügen, die im Voraus berechnet werden, um den nächsten besten Zug von einer bestimmten Schachposition zu bestimmen.
Sie wird als abhängige Variable betrachtet, da sie von der Rechenzeit beeinflusst wird, und dient eher der erklärenden Analyse als der direkten Bewertung.
Je schneller eine Position ausgewertet werden kann, desto mehr Züge im Voraus können innerhalb einer angemessenen Zeitspanne analysiert werden.

\subsection{Testaufbau}
Da die Zugtiefe direkt von der Rechenzeit abhängig ist, wurde hier der gleiche \gls{glos:JUnit}-Test verwendet, wie für die Rechenzeitberechnung (Kapitel~\ref{sec:rechenzeit}).
Die Ausführung dieser Tests ist sehr zeitaufwändig, vor allem bei grösseren Zugtiefen.
Darum wurde entschieden, dass die maximale Zugtiefe eines Algorithmus erreicht wurde, sobald die Rechenzeit 10 Sekunden überschreitet.
Diese Grenze wurde nicht technisch durchgesetzt, sondern durch das wiederholte Ausführen des Tests für alle Positionen mit unterschiedlichen Tiefen.
Sobald die Rechenzeit überzeugend 10 Sekunden überschritt, wurde diese Tiefe fest im Test hinterlegt. \\
Eine Ausnahme bildete hier die sequenzielle Implementierung des Minimax-Algorithmus ohne jegliche Optimierungen (Kapitel~\ref{subsec:minimax-sequenziell}), um einen besseren Vergleich zu ermöglichen.
Folglich wurde der Test vollständig ausgeführt, womit die genauen Rechenzeiten dokumentiert wurden. \\
Die finale Testausführung lief knapp unter 7 Tagen.

\subsection{Testresultate}
Aus den Diagrammen im Anhang (Rechenzeiten pro Zugtiefe~\ref{sec:rechenzeiten-pro-zugtiefe}) ergibt sich die folgende maximale Zugtiefe, die mit dem Testsetup innerhalb von 10 Sekunden berechnet werden kann.
Dabei wurde die Rechenzeit als Median aus 10 Messungen bestimmt.
\begin{table}[H]
    \centering
    \resizebox{\textwidth}{!}{%
        \begin{tabular}{|l|r|}
            \hline
            \rowcolor{lightgray}
            \textbf{Algorithmus} & \textbf{Maximale Tiefe} \\
            \hline
            \rowcolor[HTML]{FFFFFF}
            MiniMaxSequential & 3 \\
            \hline
            \rowcolor[HTML]{EFEFEF}
            MiniMaxSequentialWithCache & 4 \\
            \hline
            \rowcolor[HTML]{FFFFFF}
            MiniMaxParallel & 4 \\
            \hline
            \rowcolor[HTML]{EFEFEF}
            MiniMaxParallelWithCache & 4 \\
            \hline
            \rowcolor[HTML]{FFFFFF}
            MiniMaxAlphaBetaSequential & 6 \\
            \hline
            \rowcolor[HTML]{EFEFEF}
            MiniMaxAlphaBetaSequentialWithCache & 6 \\
            \hline
        \end{tabular}%
    }
    \caption{Maximale Zugtiefe innerhalb von 10 Sekunden}
    \label{tab:maximale-zugtiefe-pro-algorithmus}
\end{table}